% ============================================================
%  chapters/05_methodology.tex
% ============================================================

\section{Methodology}
\label{sec:methodology}

In this paper I estimate the effect of community-level violence exposure on
two distinct educational outcomes: current school enrollment and cumulative
years of schooling. I define educational attainment as the number of years
of school a student has completed.

\subsection{School Enrollment (\textit{attend})}
\label{sec:method_attend}

To capture the effects of violence on school enrollment, the dependent
variable is \textit{attend}---a binary indicator equal to one if the student
is currently enrolled in school and zero otherwise. The baseline estimating
equation is:

\begin{equation}
    \label{eq:attend}
    \textit{attend}_{i} = \alpha + \boldsymbol{\beta}\,\mathbf{violence}_{i}
        + \boldsymbol{\gamma}\,\mathbf{student}_{i}
        + \boldsymbol{\delta}\,\mathbf{X}_{i}
        + \varepsilon_{i}
\end{equation}

\noindent where $\mathbf{violence}_{i}$ is a vector of binary indicators for whether
the child's LGA experienced at least one UCDP-recorded violent event during
each of the following age windows: pre-birth (17 to 9 months before birth),
in utero (0 to 9 months before birth), early childhood (ages 0--3), preschool
(ages 3--6), primary school age (ages 6--12), junior secondary age (ages
12--15), and senior secondary age (ages 15--18).

$\mathbf{student}_{i}$ is a vector of individual-level student
characteristics including age, age-squared, gender, religion, an indicator
for late school start, Quranic school attendance, Tsangaya school
attendance, Islamiyya school attendance, disability indicators (seeing,
hearing, speaking, mobility, mental, other), and distance to primary,
junior secondary, and senior secondary schools. $\mathbf{X}_{i}$ contains
additional household and community-level controls: mother's and father's
living status, whether each parent resides in the household, parental
educational attainment, community wealth quintile, region indicators, and
the number of eligible school-age children in the household.

I estimate Equation~\eqref{eq:attend} by OLS with heteroskedasticity-robust
standard errors, and I also estimate an augmented version that includes LGA
fixed effects to absorb unobserved time-invariant LGA-level characteristics:

\begin{equation}
    \label{eq:attend_fe}
    \textit{attend}_{i} = \alpha + \boldsymbol{\beta}\,\mathbf{violence}_{i}
        + \boldsymbol{\gamma}\,\mathbf{student}_{i}
        + \boldsymbol{\delta}\,\mathbf{X}_{i}
        + \mu_{\ell} + \varepsilon_{i}
\end{equation}

\noindent where $\mu_{\ell}$ denotes LGA fixed effects.

\paragraph{Violence variable construction.}
The violence binary indicators in $\mathbf{violence}_{i}$ were created using
the UCDP event start dates and each child's estimated age. The household
survey provided month and year of birth for each child; I assigned a birth
date of the 15th of the reported birth month. I then calculated the
difference between each child's birth date and the start date of the first
recorded UCDP event in their LGA, if the LGA experienced any violence, to
determine the child's age at first community violence exposure. A child was
coded as not having experienced violence in a given window if born more than
17 months after the first recorded event in their LGA, since we do not
expect violence that occurred well before a child's birth to affect that
child. The 17-month pre-birth window additionally allows estimation of the
effects of a mother's prenatal living conditions on child outcomes.%
\footnote{See Appendix~\ref{app:variable_defs} for a complete description of
all constructed variables and their definitions.}

\paragraph{Late school start.}
The variable \textit{late}---a binary indicator for whether a student began
school after reaching age six---was constructed from survey questions asking
about reasons for late school entry, rather than from the age-at-entry
variable in the original dataset, which contained implausible values
(including ages younger than the student's current age).

\paragraph{Quranic, Tsangaya, and Islamiyya variables.}
These three variables were present in the original dataset but contained
missing data for a large share of observations. To address this, I assumed
that only Muslim children are likely to attend one of these school types.
Accordingly, for children who identified as Christian, Traditionalist, or
other non-Muslim religion, all three variables were set to zero. This
resolved most missing values; the remaining missing observations were
retained as missing.

\paragraph{Disability variables.}
The original disability variable was a string in which each disability
category was represented by a letter. I decomposed this string into six
separate binary indicators: \textit{seeing}, \textit{hearing},
\textit{speaking}, \textit{mobility}, \textit{mental}, and
\textit{dis\_other}.%
\footnote{The string was split into six single-character variables; a binary
indicator for each disability was then set to 1 if the corresponding
character appeared in the original string.}

\paragraph{Parental education.}
The variables on mother's and father's educational attainment contained
missing values for many observations, typically because the level of
schooling reached was recorded but the attendance indicator was missing. For
observations where the school level variable was non-zero, I imputed school
attendance as one. The variables \textit{mom\_enrol} and
\textit{dad\_enrol}---total years of schooling---were constructed by summing
years implied by the level and class reached (e.g., if a parent reached
secondary school level 2, class 3, years of enrollment = $6 + 3 = 9$). This
approach does not account for grade repetition, which the survey does not
record for parents.

\subsection{Educational Attainment (\textit{highest\_enrolled})}
\label{sec:method_attainment}

To capture the effects of violence on cumulative educational attainment, I
use the variable \textit{highest\_enrolled}, which measures the total years
of schooling completed. The estimating equation parallels
Equation~\eqref{eq:attend}:

\begin{equation}
    \label{eq:attainment}
    \textit{highest\_enrolled}_{i} = \alpha
        + \boldsymbol{\beta}\,\mathbf{violence}_{i}
        + \boldsymbol{\gamma}\,\mathbf{student}_{i}
        + \boldsymbol{\delta}\,\mathbf{X}_{i}
        + \varepsilon_{i}
\end{equation}

\noindent with an augmented version including LGA fixed effects.

\paragraph{Construction of \textit{highest\_enrolled}.}
The outcome variable was constructed to reflect the maximum years of
schooling completed by each student across three sources of information in
the dataset: (1) current class and school level for currently enrolled
students, (2) class and level from the previous school year (2008--2009),
and (3) highest class and level reached for students who had dropped out.
For currently enrolled students, the current class was assigned as the
highest class reached. For each school level and class combination, I
converted to years of schooling by summing the years required to reach that
level (e.g., a student at senior secondary school level 1, class 3 has
completed $6 + 3 + 3 = 12$ years). I then took the maximum across all three
sources to handle data entry errors, as a number of currently enrolled
students recorded a current grade lower than the previous year's grade. This
approach assumes the highest recorded grade is the true grade reached, which
may introduce upward bias for a small number of observations.
